\subsection{SEMI-SIMPLE LIE GROUPS}

PROPOSITION 26. Let \(G\) be a finite-dimensional connected real or complex Lie group. The following conditions are equivalent:

\begin{enumerate}
\def\labelenumi{(\roman{enumi})}
\item
  L(G) is semi-simple;
\item
  the radical of G is \{e\};
\item
  every normal commutative integral subgroup of \(\mathbf{G}\) is equal to \(\{e\}\) .
\end{enumerate}

Condition (ii) means that the radical of L(G) is \(\{0\}\) and hence (i) \(\Leftrightarrow\) (ii) (Chapter I, § 6, Theorem 1). The equivalence of (i) and (iii) follows from § 6, no. 6, Proposition 14.

DEFINITION 2. A connected real or complex Lie group is called semi-simple if it is finite-dimensional and satisfies the conditions of Proposition 26.

Remark 1. Let G be a finite-dimensional connected real or complex Lie group. If G is not semi-simple, G has a commutative connected Lie subgroup \(G'\) which is invariant under every continuous automorphism, such that \(G' \neq \{\varepsilon\}\) . For let n be the largest nilpotent ideal of L(G); then \(n \neq \{0\}\) and the corresponding analytic subgroup N is a Lie subgroup which is invariant under every continuous automorphism of G (no. 7, Proposition 23); the centre G' of N has the desired properties.

PROPOSITION 27. Let \(G\) be a finite-dimensional connected real or complex Lie group. The following conditions are equivalent:

\begin{enumerate}
\def\labelenumi{(\roman{enumi})}
\item
  L(G) is simple;
\item
  the only normal integral subgroups of \(\mathbf{G}\) are \(\{e\}\) and \(\mathbf{G}\) and further \(\mathbf{G}\) is not commutative.
\end{enumerate}

This follows from § 6, no. 6, Proposition 14.

DEFINITION 3. A connected real or complex Lie group is called almost simple if it is finite-dimensional and satisfies the conditions of Proposition 27.

PROPOSITION 28. Let \(\mathbf{G}\) be a simply connected real or complex Lie group. The following conditions are equivalent:

\begin{enumerate}
\def\labelenumi{(\roman{enumi})}
\item
  G is semi-simple;
\item
  G is isomorphic to the product of a finite number of almost simple groups.
\end{enumerate}

If G is a finite product of almost simple Lie groups, \(\mathrm{L}(\mathrm{G})\) is a finite product of simple Lie algebras and is hence semi-simple. If G is semi-simple, \(\mathrm{L}(\mathrm{G})\) is isomorphic to a product of simple Lie algebras \(L_{1},\ldots,L_{n}\) . Let \(G_{i}\) be a simply connected Lie group with Lie algebra \(L_{i}\) , which is therefore almost simple. Then G and \(G_{1}\times\cdots\times G_{n}\) are simply connected and have isomorphic Lie algebras and are therefore isomorphic.

Lemma 1. Let \(\mathbf{G}\) be a connected topological group, \(\mathbf{Z}\) its centre and \(\mathbf{Z}'\) a discrete subgroup of \(\mathbf{Z}\) . Then the centre of \(\mathbf{G}/\mathbf{Z}'\) is \(\mathbf{Z}/\mathbf{Z}'\) .

Let \(y\) be an element of G whose class modulo \(\mathbf{Z}'\) is a central element of

G/Z'. Let \(\phi\) be the mapping \(g \mapsto gyg^{-1}y^{-1}\) of G into G. Then \(\phi(G)\) is connected and contained in Z' and hence \(\phi(G) = \phi(\{e\}) = \{e\}\) . Therefore \(y \in Z\) .

PROPOSITION 29. Let G be a semi-simple connected real or complex Lie group.

\begin{enumerate}
\def\labelenumi{(\roman{enumi})}
\item
  G = (G, G).
\item
  The centre Z of G is discrete.
\item
  The centre of G/Z is \{e\}.
\end{enumerate}

Assertion (i) follows from the Corollary to Proposition 4, no. 2 and Chapter I, § 6, Theorem 1.

Assertion (ii) follows from § 6, no. 4, Corollary 4 to Proposition 10 and Chapter I, § 6, no. 1, Remark 2.

Assertion (iii) follows from (ii) and Lemma 1.

PROPOSITION 30. (i) Let g be a semi-simple real or complex Lie algebra. Then Int g is the identity component of Aut g.

\begin{enumerate}
\def\labelenumi{(\roman{enumi})}
\setcounter{enumi}{1}
\tightlist
\item
  Let \(\mathbf{G}\) be a semi-simple connected real or complex Lie group. The adjoint group of \(\mathbf{G}\) is the identity component of \(\operatorname{Aut} \mathrm{L}(\mathrm{G})\) . Its centre reduces to the identity element.
\end{enumerate}

Every derivation of g is inner (Chapter I, § 6, Corollary 3 to Proposition 1) and hence L(Int g) = L(Aut g), which proves (i). The first assertion of (ii) follows from (i). The second follows from Proposition 29 (iii) and § 6, no. 4, Corollary 4 (ii) to Proposition 10.

Remark 2. Let g be a complex semi-simple Lie algebra and \(g_{0}\) its underlying real Lie algebra. Then \(\operatorname{Aut}(g)\) is open in \(\operatorname{Aut}(g_{0})\) , for \(\operatorname{Int}(g_{0}) \subset \operatorname{Aut}(g)\) .

PROPOSITION 31. Let G be a simply connected finite-dimensional real or complex Lie group and R its radical. There exists a semi-simple simply connected Lie subgroup S of G such that G, as a Lie group, is the semi-direct product of S by R. If S' is a semi-simple integral subgroup of G, there exists x in the nilpotent radical of L(G) such that

\[
(\operatorname{Ad} \exp x) (\mathrm{S} ^ {\prime}) \subset \mathrm{S}.
\]

This follows from § 6, no. 6, Corollary 1 to Proposition 14 and Chapter I, § 6, Theorem 5 and Corollary 1.

Lemma 2. Let G be a group (resp. a topological group), G' a normal subgroup of G, V a finite-dimensional vector space over a commutative field k (resp. over K), \(\rho\) a linear representation (resp. a continuous linear representation) of G on V and \(\rho' = \rho|G'\) .

\begin{enumerate}
\def\labelenumi{(\roman{enumi})}
\item
  If \(\rho\) is semi-simple, \(\rho'\) is semi-simple.
\item
  If \(\rho'\) is semi-simple and every finite-dimensional linear \(k\) -representation (resp. continuous linear K-representation) of \(\mathrm{G} / \mathrm{G}'\) (resp. \(\mathrm{G} / \overline{\mathrm{G}}'\) ) is semi-simple, then \(\rho\) is semi-simple.
\end{enumerate}

Suppose that \(\rho\) is semi-simple; we prove that \(\rho'\) is semi-simple. It suffices to consider the case where \(\rho\) is simple. Let \(V'\) be a minimal non-zero sub- \(G'\) -module of \(V\) . For all \(g \in G\) , \(\rho(G') \rho(g)V' = \rho(g)\rho(G')V' = \rho(g)V'\) , in other words \(\rho(g)\mathrm{V}'\) is stable under \(\rho(\mathrm{G}')\) ; if \(\mathrm{V}''\) is a sub- \(\mathrm{G}'\) -module of \(\rho(g)\mathrm{V}'\) , then \(\rho(g)^{-1}\mathrm{V}''\) is a sub- \(\mathrm{G}'\) -module of \(\mathrm{V}'\) and hence \(\mathrm{V}''\) is equal to \(\{0\}\) or \(\rho(g)\mathrm{V}'\) . Thus, for all \(g \in \mathrm{G}\) , \(\rho(g)\mathrm{V}'\) is a simple \(\mathrm{G}'\) -module. But \(\sum_{g \in \mathrm{G}} \rho(g)\mathrm{V}'\) is a non-zero sub-G-module of \(\mathrm{V}\) , whence \(\mathrm{V} = \sum_{g \in \mathrm{G}} \rho(g)\mathrm{V}'\) . Hence \(\rho'\) is semi-simple.

Suppose that \(\rho'\) is semi-simple. Let \(W\) be a non-zero sub-G-module of \(V\) . As \(\rho'\) is semi-simple, there exists a projector \(f_0\) of \(V\) onto \(W\) which commutes with \(\rho'(G)\) . Let \(E\) be the set of \(f \in \mathcal{L}(V, V)\) which commute with \(\rho(G')\) , which map \(V\) into \(W\) and whose restriction to \(W\) is a homothety; for \(f \in E\) , let \(\alpha(f)\) denote the ratio of the homothety \(f|W\) . Then \(f_0 \in E\) and \(\alpha(f_0) = 1\) . Clearly \(\alpha\) is a linear form on \(E\) . Let \(F = \text{Ker } \alpha\) , which is a hyperplane of \(E\) . For \(f \in E\) and \(g \in G\) , we write \(\sigma(g)f = \rho(g) \circ f \circ \rho(g)^{-1}\) ; then \(\sigma(g)f\) maps \(V\) into \(W\) and its restriction to \(W\) is the homothety of ratio \(\alpha(f)\) ; if \(g' \in G'\) , then

\[
\begin{array}{r l} \sigma (g) f \circ \rho (g ^ {\prime}) & = \rho (g) \circ f \circ \rho (g) ^ {- 1} \circ \rho (g ^ {\prime}) \\ & = \rho (g) \circ f \circ \rho (g ^ {- 1} g ^ {\prime} g) \circ \rho (g ^ {- 1}) \\ & = \rho (g) \circ \rho (g ^ {- 1} g ^ {\prime} g) \circ f \circ \rho (g ^ {- 1}) \\ & = \rho (g ^ {\prime}) \circ \rho (g) \circ f \circ \rho (g ^ {- 1}) \\ & = \rho (g ^ {\prime}) \circ \sigma (g) f. \end{array}
\]

Hence \(\sigma(g)f\in\mathbf{E}\) . Therefore \(\sigma\) is a linear \(k\) -representation (resp. continuous linear K-representation) of G on E leaving F stable. Then \(\sigma(g)=\mathrm{Id}_{\mathbf{E}}\) for \(g\in\mathbf{G}'\) and hence for \(g\in\overline{\mathbf{G}'}\) in the topological case. Suppose that every finite-dimensional \(k\) -linear representation (resp. continuous K-linear representation) of G/G' (resp. G/ \(\overline{\mathbf{G}'}\) ) is semi-simple. Then there exists in E a supplement of F which is stable under G. In other words, there exists \(f\in\mathbf{E}\) such that \(\alpha(f)=1\) , which is invariant under G. Then \(f\) is a projector of V onto W, and, for \(g\in\mathbf{G}\) , \(\rho(g)\circ f\circ\rho(g^{-1})=f\) , that is \(f\) commutes with \(\rho(\mathbf{G})\) . Thus \(\rho\) is semi-simple.

THEOREM 1. Let G be a finite-dimensional real or complex Lie group, \(G_{0}\) its identity component, R its radical and r the radical of L(G); suppose that \(G/G_{0}\) is finite. Let \(\rho\) be a finite-dimensional analytic linear representation of G. The following conditions are equivalent:

\begin{enumerate}
\def\labelenumi{(\roman{enumi})}
\item
  \(\rho\) is semi-simple;
\item
  \(\rho|G_{0}\) is semi-simple;
\item
  \(\rho|R\) is semi-simple;
\item
  L(ρ) is semi-simple;
\item
  L(ρ)\textbar r is semi-simple.
\end{enumerate}

(i) \(\Leftrightarrow\) (ii) by Lemma 2 and Integration, Chapter VII, §3, Proposition 1. (ii) \(\Leftrightarrow\) (iv) and (iii) \(\Leftrightarrow\) (v) by §6, no. 5, Corollary 2 to Proposition 13. (iv) \(\Leftrightarrow\) (v) by Chapter I, §6, Theorem 4.

COROLLARY 1. Let \(\rho, \rho_{1}, \rho_{2}\) be finite-dimensional semi-simple analytic linear representations of G and n an integer \(\geqslant 0\) . Then \(\rho_{1} \otimes \rho_{2}\) , \(T^{n}\rho\) , \(S^{n}\rho\) , \(\bigwedge^{n}\rho\) (Appendix) are semi-simple.

The semi-simplicity of \(\rho_{1} \otimes \rho_{2}\) follows from Theorem 1 and Chapter I, §6, Corollary 1 to Theorem 4. The semi-simplicity of \(T^{n}_{\rho}, S^{n}_{\rho}, \bigwedge^{n}_{\rho}\) follows from the semi-simplicity of \(\rho_{1} \otimes \rho_{2}\) .

We shall see later that, if \(k\) is a commutative field of characteristic 0, \(\Gamma\) a group and \(\rho_1\) and \(\rho_2\) finite-dimensional semi-simple linear \(k\) -representations of \(\Gamma\) , then \(\rho_1 \otimes \rho_2\) is semi-simple.

COROLLARY 2. Let \(\rho\) be a finite-dimensional semi-simple analytic linear representation of \(G\) on a vector space \(V\) , \(S\) the symmetric algebra of \(V\) and \(S^{\mathfrak{G}}\) the subalgebra of \(S\) consisting of the elements invariant under \((\mathsf{S}(\rho))(G)\) . Then \(S^{\mathfrak{G}}\) is a finitely generated algebra.

This follows from Theorem 1, Chapter I, §6, Theorem 6 (a) and Commutative Algebra, Chapter V, §1, Theorem 2.

COROLLARY 3. Let G be a real or complex Lie group and \(G_{0}\) its identity component. Suppose that \(G_{0}\) is semi-simple and that \(G/G_{0}\) is finite. Then every finite-dimensional analytic linear representation of G is semi-simple.

PROPOSITION 32. Let G be a finite-dimensional connected real Lie group. Suppose that L(G) is reductive. The following conditions are equivalent:

\begin{enumerate}
\def\labelenumi{(\roman{enumi})}
\item
  \(\mathbf{G} / \overline{\mathbf{D}^1}\mathbf{G}\) is compact;
\item
  (resp. (ii')) every finite-dimensional analytic linear representation of \(G\) on a complex (resp. real) vector space is semi-simple.
\item
  \(\Rightarrow\) (ii'): Suppose that \(G / \overline{D}^1 G\) is compact. Then every continuous linear representation of \(G / \overline{D}^1 G\) on a finite-dimensional real vector space is semi-simple (Integration, Chapter VII, §3, Proposition 1). Let \(\rho\) be a finite-dimensional analytic linear representation of \(G\) on a real vector space. Then \(\rho |D^1 G\) is analytic, \(D^1 G\) is semi-simple (Chapter I, §6, Proposition 5) and hence \(\rho |D^1 G\) is semi-simple (Corollary 3 to Theorem 1). Hence \(\rho\) is semi-simple (Lemma 2).
\end{enumerate}

It can be seen similarly that (i) \(\Rightarrow\) (ii).

\begin{enumerate}
\def\labelenumi{(\roman{enumi})}
\tightlist
\item
  \(\Rightarrow\) (ii): Suppose that \(G / \overline{D^1} G\) is not compact and hence isomorphic to a group of the form \(\mathbf{R}^p\times \mathbf{T}^q\) with \(p > 0\) (§6, no. 4, Proposition 11 (ii)). Then there exists a surjective morphism of \(G / \overline{D^1} G\) onto \(\mathbf{R}\) and hence a surjective morphism of \(G\) onto \(\mathbf{R}\) . The mapping
\end{enumerate}

\[
g \mapsto \sigma (g) = \left( \begin{array}{c c} 1 & 0 \\ \rho (g) & 1 \end{array} \right)
\]

is an analytic linear representation of G on \(R^{2}\) which is not semi-simple, for the only 1-dimensional vector subspace of \(\mathbf{R}^2\) which is stable under \(\sigma(G)\) is \(\mathbf{R}(0,1)\) .

It can be seen similarly that (ii) \(\Rightarrow\) (i).

PROPOSITION 33. Let \(G\) be a finite-dimensional complex Lie group whose number of connected components is finite, \(\rho\) a finite-dimensional analytic linear representation of \(G\) and \(G'\) an integral subgroup of the real Lie group \(G\) such that \(L(G')\) generates \(L(G)\) over \(C\) . Then, for \(\rho\) to be semi-simple, it is necessary and sufficient that \(\rho |G'\) be semi-simple.

Let \(\rho' = \rho|G'\) . For \(\rho\) (resp. \(\rho'\) ) to be semi-simple, it is necessary and sufficient that \(L(\rho)\) (resp. \(L(\rho')\) ) be semi-simple (Theorem 1). Let \(V\) be the space of \(\rho\) . For a vector subspace of \(V\) to be stable under \(L(\rho)(L(G))\) , it is necessary and sufficient that it be stable under \(L(\rho')(L(G'))\) . Hence the proposition.

